\documentclass[11pt,a4paper]{article}
% Shared preamble for the English probability and statistics documents.
% Based on the original course preamble. Compile with pdfLaTeX.
\usepackage[top=1.5cm,bottom=2cm,left=2cm,right=2cm]{geometry}
\usepackage[T1]{fontenc}
\usepackage{lmodern}
\usepackage[utf8]{inputenc}
\usepackage[english]{babel}

\usepackage{amsmath}
\usepackage{amssymb}
\usepackage{amsthm}
\usepackage{mathrsfs}
\usepackage{mathtools}
\usepackage{etoolbox}
\usepackage{tikz}
\usetikzlibrary{shapes.multipart,arrows.meta}
\usepackage{xcolor}
\usepackage{enumerate}
\usepackage[shortlabels]{enumitem}
\usepackage{hyperref}
\usepackage{version}
\usepackage{bbold}
\usepackage{fontawesome5}
\usepackage{fancyhdr}
\lfoot{}
\usepackage{natbib}
\usepackage{graphicx}
\usepackage{pdfpages}
\usepackage{authblk}
    
\usepackage{blindtext}
\usepackage{multicol}
\usepackage{collcell}
\usepackage{wrapfig}
\usepackage{blkarray}
\usepackage{cancel}
\usepackage{appendix}
\usepackage{float}
\makeatletter
\input{tex/assets/slashbox.sty}
\makeatother
\hypersetup{unicode,colorlinks,citecolor=red,filecolor=black,
  linkcolor=red,urlcolor=cyan}
\setlength{\columnsep}{1cm}
\pagenumbering{gobble}

\def\numtext#1{%
  \ifcase#1\or one\or two\or three\or four\or five\or six\or seven\or eight\or nine\or ten\or
  eleven\or twelve\or thirteen\or fourteen\or fifteen\or sixteen\or seventeen\or eighteen\or nineteen\or twenty\or Lots\fi}

\newcommand{\Byte}[1]{%
  \def\nodecontents{}%
  \foreach[count=\l] \k in {#1} {\xappto{\nodecontents}{\noexpand\nodepart{\numtext{\l}}\k}}%
  \node[rectangle split,rectangle split parts=8,rectangle split horizontal,draw] at (2,2){\nodecontents};}

\newcommand{\D}{{\mathrm{d}}}
\newcommand{\cE}{{\mathcal{E}}}
\newcommand{\cG}{{\mathcal{G}}}
\newcommand{\cX}{{\mathcal{X}}}
\newcommand{\cY}{{\mathcal{Y}}}
\newcommand{\cB}{{\mathcal{B}}}
\newcommand{\cF}{{\mathcal{F}}}
\newcommand{\cO}{{\mathcal{O}}}
\newcommand{\cS}{{\mathcal{S}}}
\newcommand{\cT}{{\mathcal{T}}}
\newcommand{\given}{\, \big| \,}
\newcommand{\cL}{{\mathscr{L}}}
\newcommand{\cD}{{\mathcal{D}}}
\newcommand{\cA}{{\mathscr{A}}}
\newcommand{\cP}{{\mathscr{P}}}
\newcommand{\cR}{{\mathscr{R}}}
\newcommand{\A}{{\mathcal A}}
\newcommand{\one}{\mathbb{1}}
\newcommand{\mathbbm}[1]{\mathbb{#1}}

\providecommand{\C}{}
\renewcommand{\C}{\mathbb C}
\newcommand{\R}{\mathbb R}
\newcommand{\Q}{\mathbb Q}
\newcommand{\N}{\mathbb N}
\newcommand{\T}{\mathbb T}
\newcommand{\F}{\mathbb F}
\renewcommand{\P}{\mathbb P}
\newcommand{\Z}{\mathbb Z}
\newcommand{\E}{\mathbb {E}}
\newcommand{\Pbb}[1]{\P\left( #1\right)}
\newcommand{\Ebb}[1]{\E\left[ #1\right]}

\newcommand{\lb}{\left (}
\newcommand{\rb}{\right )}
\newcommand{\lbb}{\left [}
\newcommand{\rbb}{\right ]}
\newcommand{\labs}{\left |}
\newcommand{\rabs}{\right |}
\newcommand{\lbrb}[1]{\lb #1 \rb}
\newcommand{\lbbrbb}[1]{\lbb#1\rbb}
\newcommand{\labsrabs}[1]{\labs#1\rabs}
\newcommand{\lbbrb}[1]{\lbb#1\rb}
\newcommand{\lbrbb}[1]{\lb#1\rbb}
\newcommand{\langlerangle}[1]{\langle#1\rangle}
\newcommand{\lbcurly}{\{}
\newcommand{\rbcurly}{\}}
\newcommand{\curly}[1]{\lbcurly#1\rbcurly}
\newcommand{\abs}{\labsrabs}
\newcommand{\ind}[1]{1_{#1}}
\newcommand{\indep}{\perp \!\!\! \perp}
\renewcommand{\d}{\, \mathrm{d}}

\newtheorem{theorem}{Theorem}
\newtheorem{problemCnt}{Lemma}
\theoremstyle{definition}
\newtheorem{problem}[problemCnt]{Problem}
\newtheorem{problembook}[problemCnt]{\faBook \,  Problem}
\newtheorem{problempencil}[problemCnt]{\faPencil* \,  Problem}
\newtheorem{problemstar}[problemCnt]{\faStar\,  Problem}
\newtheorem{problemstarpencil}[problemCnt]{\faStar \, \faPencil* \,   Problem}
\newtheorem{probleminteresting}[problemCnt]{\faSmile[regular] \,   Problem}
\newtheorem{solution}{Solution}

\newcommand{\fwcell}[1]{\makebox[\arraystretch\normalbaselineskip]{$#1$}}
\newenvironment{enumeratei}{\begin{enumerate}[(i)]}{\end{enumerate}}
\newenvironment{enumeratea}{\begin{enumerate}[(a)]}{\end{enumerate}}
\newenvironment{enumeraten}{\begin{enumerate}[\upshape 1.]}{\end{enumerate}}

% Compatibility helpers needed by a few files in the combined archive.
\AtBeginDocument{\noindent{\footnotesize English translation prepared with ChatGPT (OpenAI, GPT-6), 26 October 2026, for translation of the Bulgarian text into English and LaTeX formatting.
The Bulgarian originals remain the reference; translation may contain errors.}\par\smallskip}
\newcommand{\MissingGraphic}[1]{%
  \begin{center}\fbox{\parbox{0.82\linewidth}{\centering Missing source image: \texttt{#1}}}\end{center}}



\pagestyle{fancy}\fancyhf{}\lhead{Probability and Statistics: Resit solutions}\rhead{17 August 2022}
\setlength{\headheight}{15pt}\pagenumbering{arabic}\cfoot{\thepage}
\begin{document}
\noindent\emph{English transcription of the original handwritten solutions. Apparent source inconsistencies are identified in notes.}

\section*{Problem 1}
Filling the first square takes $\operatorname{Ge}(1)$ trials; filling a second distinct square after the first takes $\operatorname{Ge}(8/9)$, and so on. Hence
\[
\E[\text{time}]=\E[\operatorname{Ge}(1)+\cdots+\operatorname{Ge}(1/9)]
=9\left(1+\frac12+\cdots+\frac19\right)\approx25.46.
\]
Let $T_2$ be the first time at which one square contains two darts. Split according to the square where this first repetition occurs. By symmetry,
\begin{align*}
\E T_2
&=9\E[T_2\mathbf1_{\{\text{first repetition occurs in a specified square}\}}]\\
&=9\left[2\left(\frac19\right)^2
+3\binom81\,2!\left(\frac19\right)^3
+4\binom82\,3!\left(\frac19\right)^4+\cdots\right].
\end{align*}
In compact form,
\[
\E T_2=9\sum_{k=0}^8(k+2)!\left(\frac19\right)^{k+2}\binom8k
\approx4.46.
\]
For an $n\times n$ grid, the answer to the first question is
\[
n^2\left(1+\cdots+\frac1{n^2}\right)\sim n^2\ln(n^2)=2n^2\ln n.
\]
For the second question, the handwritten source gives
\[
n^2\sum_{k=0}^{n-1}(k+2)!\left(\frac1{n^2}\right)^{k+2}\binom{n-1}k.\tag{Note 1}
\]

\section*{Problem 2}
\paragraph{1.} Since $(X_1,\ldots,X_n)$ is a permutation,
\[
S=X_1+\cdots+X_n=1+\cdots+n=\frac{n(n+1)}2.
\]
Therefore $DS=0$ and $\E S=n(n+1)/2$.
\paragraph{2.}
\begin{align*}
D(X+Y)&=\E[(X+Y)^2]-(\E X+\E Y)^2\\
&=\E X^2-(\E X)^2+\E Y^2-(\E Y)^2
+2(\E[XY]-\E X\E Y).
\end{align*}
\paragraph{3.} By symmetry,
\[
\E S=\E(X_1+\cdots+X_n)=n\E X_1,
\qquad\E X_i=\frac{n+1}2.
\]
Similarly,
\[
\E(X_1^2+\cdots+X_n^2)=\E(1^2+\cdots+n^2)
=\frac{n(n+1)(2n+1)}6=n\E X_1^2,
\]
so
\[
\E X_i^2=\frac{(n+1)(2n+1)}6,
\]
\[
DX_i=\E X_i^2-(\E X_i)^2
=\frac{n+1}2\left(\frac{2n+1}3-\frac{n+1}2\right)=\frac{n^2-1}{12}.
\]
\paragraph{4.}
\begin{align*}
0=DS&=nDX_1+\binom n2\,2\operatorname{Cov}(X_1,X_2)\\
&=n\frac{n^2-1}{12}+n(n-1)\operatorname{Cov}(X_1,X_2).
\end{align*}
Hence
\[
\operatorname{Cov}(X_i,X_j)=
\begin{cases}-(n+1)/12,&i\ne j,\\(n^2-1)/12,&i=j.\end{cases}
\]

\section*{Problem 3}
\paragraph{1.} Independence gives $D(M_1+\cdots+M_n)=nDM_1$. Let $M=M_1=\max(U_1,V_1)$. For $t\in[0,1]$,
\[
\P(M\le t)=\P(U_1\le t)\P(V_1\le t)=t^2,
\qquad f_M(t)=2t\mathbf1_{[0,1]}(t).
\]
Therefore
\[
\E M=\int_0^12t\cdot t\d t=\frac23,
\qquad\E M^2=\int_0^1t^2\cdot2t\d t=\frac12,
\]
\[
DM=\frac12-\frac49=\frac1{18},\qquad
D(M_1+\cdots+M_n)=\frac n{18}.
\]
\paragraph{2.} The source writes the CLT approximation as
\[
M_1+\cdots+M_n\approx N\left(n\cdot\frac23,\frac1{18}\right).\tag{Note 2}
\]

\section*{Problem 4}
\paragraph{1.} The density normalization requires
\[
\int_1^\infty\frac c{x^4}\d x=1,
\qquad\left[-\frac c3x^{-3}\right]_1^\infty=\frac c3,
\qquad c=3.
\]
Next,
\[
\E X_1=\int_1^\infty\frac c{x^3}\d x
=\left[-\frac c2x^{-2}\right]_1^\infty=\frac c2=\frac32,
\]
\[
\E X_1^2=\int_1^\infty\frac c{x^2}\d x=c=3,
\qquad DX_1=3-\frac94=\frac34.
\]
Since $DX_1<\infty$, the central limit theorem gives
\[
\frac{S_n-n\cdot\frac32}{\sqrt{3n/4}}\approx N(0,1),
\qquad a=\frac{3n}2,\quad b=\frac{\sqrt{3n}}2.
\]
\paragraph{2.} Let $S=\sum_{i=1}^{40}X_i$. Here $\E S=DS=400$. Using the CLT with $40>30$,
\begin{align*}
99\%&\le\P(S>a)=\P\left(\frac{S-400}{20}>\frac{a-400}{20}\right)\\
&\approx\P\left(N(0,1)>\frac{a-400}{20}\right).
\end{align*}
From the normal table, $(a-400)/20\approx-2.33$, so
\[
a=400-46.6=353.4.
\]

\section*{Transcription notes}
\begin{enumerate}
\item For an $n\times n$ grid, the compact expression should use the $n^2-1$ other squares, so both occurrences of $n-1$ in this source formula require review.
\item The preceding variance calculation gives $n/18$. The normal approximation in the original omits the factor $n$ in its variance parameter.
\end{enumerate}
\end{document}
